\scanpage{37}\mspage{30}
% archive sheet number: 36
est une équivalence, mais non pas que le foncteur localisation
\[
 M_{cf}\unsure{-|\longrightarrow}M_{cf}W_{cf}^{-1}
\]
est un \deleted{\unsure{simple}} passage \unsure{aux|au} quotient sur les flèches, i.e. que si $X,Y$ sont dans $M_{cf}$, alors
\diagram{\begin{tikzcd}[ampersand replacement=\&,row sep=small]
\text{(a)}\qquad\operatorname{Hom}_{M}(X,Y)\arrow[r] \& \operatorname{Hom}_{W}(X,Y)\arrow[d,equal,"\mathrm{def}"]\\
\&\operatorname{Hom}_{\operatorname{Ho}_{W}}(X,Y)
\end{tikzcd}}
soit surjectif. Ceci est assuré par l'axiome d'exactitude $\mathrm M2$ bis
\[
 W=TF\circ TC.
\]

\deleted{\unsure{Ainsi}} Avec ce seul axiome, on n'a pas de caractérisation sympa de la relation d'équivalence dans $\operatorname{Hom}_{M}(X,Y)$ définie par la flèche (a) (pour $X,Y$ dans $M_{cf}$).

% slanted
\margin{Il semblerait que la relation de Q-homotopie à droite \deleted{\unsure{est}} est \deleted{\unsure{une}} \deleted{\unsure{équiv.}} une relation d'équivalence, pour $\unsure{X,Y}$ de $\unsure{M_{cf}|M_{f}}$, et cela pour la \unsure{seule}  str. de Quillen faible exacte à droite.}
Si on a exactitude à gauche
\[
 TC=C\cap W\qquad(\Longleftrightarrow\quad TC=C\cap W^{\unsure{\mathrm{univ}|\mathrm{uni}}}),
\]
alors \unsure{cette} relation d'équivalence est \deleted{\unsure{la}} engendrée \inserted{par la relation de} \uline{Q-homotopie à droite} (\unsure{impliquée} par la Q-homotopie à gauche), et la catégorie quotient
% slanted
\margin{\uline{NB} D'après \unsure{Baues|Barnes}, il semblerait que sur $M_{\unsure{cf|f}}$, cette relation d'homotopie à droite équivaut à la relation d'homotopie à gauche et que celle-ci \unsure{ne} dépend pas du ``cylindre'' choisi.}
\[
 \pi_{d}M_{cf}=\pi M_{cf}\approx\operatorname{Ho}_{W}\deleted{\unsure{M_{cf}}}
\]
s'envoie dans $\pi_{d}M_{c}$
\[
 \deleted{\unsure{\pi_{d}M_{c}}}\ \pi M_{cf}\longrightarrow\pi_{d}M_{c}
\]
et ce foncteur peut être interprété \unsure{comme} (\unsure{ce dernier|ceci}
